\documentclass[11pt]{article}

\usepackage{amsmath, amssymb}
\usepackage{geometry}
\usepackage{booktabs}
\usepackage{longtable}
\usepackage{hyperref}
\usepackage{xcolor}
\usepackage{enumitem}
\usepackage{url}

\geometry{margin=1in}

\DeclareUrlCommand\code{\urlstyle{tt}}
\setlength{\emergencystretch}{2em}
\newcommand{\status}[1]{\textcolor{gray}{\textsc{status:} #1}}

\title{Future research directions\\
\large Options after ``Driving in the wrong direction? Modelling policy
mixes for EV adoption''}
\author{}
\date{}

\begin{document}
\maketitle

\begin{abstract}
\noindent
Fourteen candidate directions for extending the agent-based model of the
ICEV--EV transition, grouped into five themes, each assessed against three
questions: which limitation of the published paper does it address, what
already exists in the repository, and what does it cost --- in particular,
whether it forces a re-run of the simulation-based calibration. The
concluding sections give a summary table and a recommended sequence. The
central practical finding is that the directions differ in cost by more than
an order of magnitude, and that the two cheapest are also the two that lift
limitations the published paper states explicitly.
\end{abstract}

\tableofcontents

\section{Purpose and how to read this}

The published model is calibrated on California 2001--2023 and used to
evaluate five market-based instruments, individually and in pairs, over
2024--2035 with a projection to 2050. It is a working, calibrated,
documented artefact. Every direction below is therefore an
\emph{extension}, not a rebuild, and the relevant question for each is not
``is this interesting'' --- they all are --- but ``what does it cost, what
does it buy, and does it break the calibration''.

Three cross-cutting facts govern everything that follows.

\begin{description}[leftmargin=1.5em]
\item[The calibration is the bottleneck.] Only four parameters are
  estimated by simulation-based inference: $a_\chi$, $b_\chi$
  (the innovativeness-threshold Beta distribution), $\kappa$ (consumer
  intensity of choice) and $\delta$ (fuel-efficiency depreciation). See
  \code{package/calibration/NN_multi_round_calibration_multi_gen.py}.
  Any extension that changes agent behaviour \emph{during} 2001--2023
  invalidates that posterior and forces a re-run. Extensions confined to the
  2024--2035 policy period do not. This single distinction separates the
  cheap directions from the expensive ones and is the axis the summary table
  in Section~\ref{sec:table} is sorted on.

\item[No firm behavioural parameter is estimated.] $\lambda$ (research
  intensity of choice, $10^{-3}$), $M^+$ (memory, 30), $Y^+$ (variant cap,
  10), $\theta$ and $\iota$ (activation probabilities) are all set by
  assumption, as Table~4 of the manuscript records. The firm side is
  disciplined only indirectly, through the price, HHI and car-age target
  ranges of Table~1. This matters because it means the supply-side
  extensions (Themes A, C, D) are not competing with estimated parameters
  --- they are replacing one asserted specification with another, and the
  case for each has to be made on economic grounds rather than fit.

\item[Several directions are already partly built.] Forward-looking
  expectations and three command-and-control bans are implemented in the
  current model code behind default-off flags. A Gaussian-process surrogate
  over the full five-instrument policy space, and a ZIP-level spatial
  calibration pipeline with a 28-dimensional summary statistic, both exist
  in git history and were removed from the publication snapshot. Each
  section below records this under \status{}.
\end{description}

\section{What the paper commits us to}
\label{sec:commitments}

The conclusions section of the manuscript names five future directions of
its own. Mapping them onto the numbered ideas is the cheapest available
source of referee-proof motivation, because a limitation the paper itself
raises does not need to be argued for.

\begin{longtable}{@{}p{0.46\textwidth}p{0.46\textwidth}@{}}
\toprule
\textbf{Stated in the paper} & \textbf{Corresponding direction} \\
\midrule
\endhead
``A key shortcoming \dots is the lack of financial and physical
constraints \dots Future research should incorporate capital costs of
switching production and a supply model with an explicit representation of
production planning and inventory management.''
& Items 5 (financial constraints), 6 (switching and R\&D costs),
12 (batch/vintage production). Theme C. \\
\addlinespace
``Incorporating more sophisticated forecasting techniques and factoring in
expected regulation can crucially modify investment and purchase decisions.
By doing so, policies such as standards and regulations could be assessed.
Additionally, this would enable the study of more realistic time-varying
carbon price implementations such as the EU-ETS.''
& Items 1 (forward-looking agents), 2 (command and control), 4 (long-horizon
firm profit), 7 (MARL firms). Theme A. Also the time-varying half of
item 8. \\
\addlinespace
``The role of pro-environmental attitudes or conspicuous consumption \dots
could be better represented by expanding \dots beyond a threshold model to
enrich car users with social norms and behavioural control such as in the
Theory of Planned Behaviour \dots This would facilitate the inclusion of
policies such as information provision, media or marketing.''
& \textbf{Not on the current list.} See Section~\ref{sec:missing}. \\
\addlinespace
``[T]ranslating these results into concrete recommendations would benefit
from moving beyond calibration to validation. Achieving this requires a
larger car user population to reduce noise \dots enabling time-series data
on EV fleet size or sales to be split into training, validation, and test
sets.''
& Items 9 (spatial), 10 (population synthesis). Theme E. \\
\addlinespace
``[W]e restrict our analysis to policy pairs rather than exhaustive
combinations to maintain tractability'' (Section~2) and policies ``are
applied from January 2024 at their full intensity and are then constant
over the policy period'' (Section~4.1).
& Item 8 (multi-policy optimisation). Theme B. \\
\bottomrule
\end{longtable}

\section{The five themes}

\begin{description}[leftmargin=1.5em]
\item[Theme A --- Expectations, anticipation and strategic firms.]
  Items 1, 2, 4, 7. All four attack the same sentence in Section~4.1 of the
  manuscript: ``we disregard the role of expectations regarding policy
  stringency \dots with agents assuming that policies will be permanent.''
  They form a ladder of increasing cost and decreasing certainty, and the
  rungs are not independent: item 2 is only meaningful given item 1, and
  item 7 delivers little that items 1 and 4 do not deliver more cheaply.

\item[Theme B --- Policy-space search and the planner's problem.]
  Item 8. Does not change the model at all; changes how the policy space is
  searched. The only theme with no recalibration exposure.

\item[Theme C --- Supply-side economics: constraints, costs, capacity.]
  Items 5, 6, 12. The firm side is where the model is most stylised, and
  this is the paper's own first-named limitation.

\item[Theme D --- Technology representation.]
  Items 3, 13, 14. How the NK landscape, battery technology and production
  carbon intensity are specified. Infrastructure work with one exception
  (item 14) that touches a headline result cheaply.

\item[Theme E --- Spatial grounding and empirical identification.]
  Items 9, 10, 11. The most expensive theme and the one with the strongest
  claim on the model's scientific credibility, because it is the only theme
  that can defend the paper's \emph{core mechanism} against its most
  obvious critique.
\end{description}

\section{Theme A: expectations, anticipation and strategic firms}

\subsection{Item 1 --- Forward-looking agents}
\label{sec:item1}

\status{implemented, documented, default off.}
See \code{docs/forward_looking/forward_looking_expectations.tex} for the
full derivation. The flag \code{forward_looking_expectations} replaces the
naive closed-form geometric series
\begin{equation}
  \mathrm{LCC}_{\mathrm{extra}} = D_i\,
  \frac{(1+r)(1-\delta)(c + \gamma_i e)}{\omega\,(r - \delta - r\delta)}
  \label{eq:naive}
\end{equation}
--- which assumes today's fuel cost and emissions intensity persist forever
--- with a backward recursion $\mathrm{Index}(t) = z_t + q\,\mathrm{Index}(t+1)$
over the model's actual known future path. It is applied consistently to
consumers (\code{socialNetworkUsers.py}) and firms
(\code{firm.py:_lifecycle_cost_term_firm}), and the carbon-price
component is extended past the simulated horizon using
\code{calculate_price_at_time()} so that a forward-looking agent
correctly believes a temporary carbon tax is temporary rather than
permanent at its peak.

\paragraph{Why it matters.} The manuscript's own Section~2.2 flags the
assumption --- ``car users assume the future prices and emissions of
gasoline and electricity to be the same as in the present period'' --- and
justifies it by citing Anderson et al.\ (2013). That justification is
reasonable for \emph{gasoline price} expectations, which is what Anderson
et al.\ study. It is much weaker for a pre-announced policy ramp, which is
public information rather than a forecasting problem.

\paragraph{Effect size is already known to be large.} The
\code{car_ban} smoke test recorded in the forward-looking note found that
because $q = 1/((1+r)(1-\delta)) \approx 0.9977$, the effective planning
horizon is $1/(1-q) \approx 433$ months. A driving ban six to eleven years
out is therefore already 70--85\% priced in from the start of the future
period. Forward-looking EV fleet share reached 0.97 six months
\emph{before} a 2030 ban took effect, against 0.39 for naive agents facing
the identical ban. Anticipation is not a marginal correction here; it is
the dominant channel.

\paragraph{Cost.} One SBI re-run. The flag changes calibration-period
behaviour too, because gasoline prices, electricity prices and grid
decarbonisation all vary over 2001--2023 even though the carbon price is
forced to zero there. It also changes how the calibrated $\delta$ enters
utility, so $\delta$ in particular should be expected to move.

\subsection{Item 2 --- Command-and-control policies}

\status{three bans implemented in \code{controller.py}; experiment packages
removed from the snapshot.}
\code{ICE_driving_ban_time} / \code{ICE_driving_ban_penalty} (a cost
shock fed through the same discounting machinery as item 1, so anticipation
emerges from the existing utility function with no separate
``anticipation lead'' parameter), \code{ICE_sales_ban_time} (a hard
constraint on firms' choice sets, which also flushes pre-ban ICE stock from
\code{cars_on_sale}) and \code{ICE_research_ban_time} (no ICE
neighbouring technologies generated at all) all exist and all default to
\code{None}. The driving-ban experiment lived in \code{package/car_ban} at
commit \code{366d91a} and is recoverable from git.

\paragraph{Why it matters most for this paper specifically.} California is
the case study \emph{and} the paradigmatic command-and-control jurisdiction
--- the manuscript says so in Section~3, cites the ZEV mandate, and then
explicitly declines to model it, on the grounds that ``regulatory measures
are out of the scope of the present research''. That is a defensible
scoping decision for one paper and an obvious opening for the next. A paper
that puts the ZEV mandate back in, and compares it against the
market-based instruments already characterised, has a ready-made
motivation, a ready-made baseline, and an audience that is already asking
the question.

\paragraph{Cost.} Essentially free \emph{given item 1}. Bans are configured
as absolute timesteps after the burn-in and can be confined to the future
period, so they do not touch the calibration. But a ban that agents cannot
anticipate is just a cost shock arriving on a known date, which is not what
a ZEV mandate is. Item 2 without item 1 answers the wrong question.

\subsection{Item 4 --- Firms chasing long-horizon profit}

\status{not implemented. The cheapest unbuilt item on the list.}

The manuscript's Section~2.3 says manufacturers ``select the car designs to
put on sale based on their profit expectations throughout the lifecycle of
the technology''. The implementation is narrower than that: Equation~8's
expected profit, as computed in
\code{firm.py:calc_predicted_profit_segments_research}, is a
\emph{single-period} quantity --- marginal profit $\times$ segment
population $\times$ expected in-segment share, all evaluated at this
month's prices, this month's subsidy scalar and this month's segment sizes.
A firm choosing a research direction today receives a technology that will
sell for years, and scores it on one month of profit.

\paragraph{The fix needs no learning.} Replace the one-period profit in
Equation~12's softmax with a discounted profit stream over the product's
expected commercial lifetime, using the discount rate $r$ the model already
has. Two sub-channels are worth separating because they can be implemented
and reported independently:
\begin{enumerate}[leftmargin=1.5em]
\item \textbf{Policy instruments on the firm's own P\&L.} The production
  subsidy and the new-EV rebate enter \code{firm.next_step()} as scalars.
  A firm weighing an R\&D direction should weight them over the payoff
  window, not this month. This is the same backward recursion item 1
  already implements, applied to the subsidy vectors.
\item \textbf{Segment growth.} $|I_{s,t}|$ and $W_{s,t}$ are today's
  segment sizes and competitive state. But EV consideration spreads
  endogenously as $\chi$ thresholds are crossed --- which is precisely the
  mechanism policy is trying to trigger --- and firms are blind to it. This
  sub-channel needs a model of segment evolution, however crude, and is
  where the honest ceiling on a non-learning implementation sits.
\end{enumerate}

\paragraph{A predicted tension worth stating in advance.} Two opposing
forces act here. Forward-looking firms that see the carbon ramp coming
should pivot to EVs \emph{earlier}, lowering the policy intensity needed to
hit 95\%. But Figure~S6 panel~G reports that EV uptake \emph{increases}
when research choices are more random, because policy incentives favouring
EVs outweigh the ICE research advantage accumulated during burn-in --- so
making firms sharper should push the other way. Against that, the Sobol
indices (Figures~S7, S8) attribute almost no global variance in EV adoption
to $\lambda$. The reasonable reading is that the \emph{level} of firm
rationality is not first-order in the current specification, but whether
firms anticipate \emph{policy} may well be. That asymmetry is an argument
for prioritising items 1 and 4 over item 7.

\subsection{Item 7 --- Strategic, pre-emptive firms via MARL}

\status{not implemented. Highest risk on the list.}

Replace the myopic softmax over NK neighbours with a learned policy
$\pi(a \mid s)$ trained to maximise discounted profit. Four objections, in
descending order of seriousness:

\begin{enumerate}[leftmargin=1.5em]
\item \textbf{It breaks the calibration chain.} SBI draws $\theta$ and
  re-simulates thousands of times. If each draw needs trained firms, the
  nested loop is infeasible. There are two ways out: freeze one policy
  trained at the posterior mode and hold it across draws (defensible, since
  $\lambda$ is held fixed across draws today), or \emph{amortise} --- put
  $\theta$ in the firm's observation vector and sample it from the SBI prior
  during training, so one trained policy is valid at any $\theta$. The
  latter is the technically right answer and has a pleasing symmetry with
  what NPE already does: train once over the prior, valid everywhere.
\item \textbf{It trades one interpretable knob for several
  uninterpretable ones.} $\lambda$ has an economic reading (exploration
  versus exploitation). Learning rate, network width, training length and
  exploration schedule do not, and they move outcomes. Any result has to
  be shown invariant to them, or reported as a range. Budget for this; it
  is the part referees will find.
\item \textbf{Multi-agent learning is non-stationary.} Ten co-learning
  firms have no convergence guarantee, and outcomes become sensitive to
  algorithmic detail --- a poor property for a model whose output is a
  policy recommendation. Best-response iteration (fix firms, train planner;
  fix planner, train firms; repeat and check for a fixed point) is far more
  stable and more reportable than simultaneous learning.
\item \textbf{Most of its value arrives more cheaply via items 1 and 4.}
  Against a \emph{pre-announced deterministic} policy path --- which is what
  \code{calculate_price_at_time()} produces --- anticipation requires no
  learning at all. A firm can read the schedule. MARL becomes genuinely
  necessary only once the planner itself is adaptive (the RL half of
  item 8), because then the policy path depends on the state, which depends
  on what firms do, and no closed form exists.
\end{enumerate}

\paragraph{If pursued, keep it minimal.} Do not learn pricing --- 
\code{calc_optimal_price_cars} is a closed-form optimum (Equation~10,
via the Lambert $W$ function) and learning can only degrade it. Replace
only \code{select_car_lambda_research}, or even just let the agent set a
scalar EV/ICE direction bias that shifts the profit scores. Share one
policy network across all ten firms, which avoids ten-way non-stationarity
and is defensible as an industry-standard heuristic.

\paragraph{The result that would justify it.} An adaptive regulator that
visibly raises subsidies when uptake stalls gives firms an incentive to
slow-roll EV investment. If the flexibility gain from an adaptive planner
partly evaporates once firms learn to exploit it, that is a
commitment-versus-discretion finding, and it is a real contribution rather
than a methods demonstration. But it presupposes item 8.

\section{Theme B: policy-space search and the planner's problem}

\subsection{Item 8 --- Multi-policy optimisation}
\label{sec:item8}

\status{Gaussian-process surrogate fully prototyped at commit
\code{c92bff8}; removed from the publication snapshot.}

The recoverable \code{package/surrogate/} contained:
\begin{itemize}[leftmargin=1.5em]
\item \code{sampling.py} --- Latin hypercube design over the full
  five-dimensional policy space, evaluating four outputs
  (\code{ev_uptake}, \code{utility}, \code{emissions}, \code{net_cost}).
\item \code{surrogate.py} --- a Gaussian process with a Mat\'ern 5/2 kernel
  plus a \code{WhiteKernel} nugget to absorb seed-to-seed variance, with
  leave-one-out cross-validation and prediction-interval coverage
  diagnostics. (The term in the original idea list, ``Gaussian mixture
  model'', is worth correcting: a GMN/mixture density network is what
  \code{sbi} uses for the \emph{posterior}; what is wanted here is a
  Gaussian \emph{process} emulator of the policy response surface. Both
  already exist in this project, doing different jobs.)
\item \code{optimisation.py} --- active Bayesian optimisation with a
  constrained acquisition function,
  $\mathrm{EI}(x) \times \Pr(\text{EV uptake} \in [\underline{u},
  \overline{u}] \mid x)$, which reduces to pure feasibility-seeking before
  the first feasible point is found; then Pareto-front extraction over
  $\sim$170 evaluated policies via randomly weighted scalarisations.
\end{itemize}

\paragraph{Why this is the strongest near-term option.} It lifts two
limitations that the paper states are computational rather than scientific.
First, pairs: ``we restrict our analysis to policy pairs rather than
exhaustive combinations to maintain tractability''. The paper's own finding
of perfect substitutability between adoption and production subsidies
(Section~4.2.2, the horizontal line in Figure~4) is a two-dimensional slice
of a surface a surrogate can map properly. Second, constant intensities:
policies are ``applied from January 2024 at their full intensity and are
then constant''. A surrogate over instrument \emph{and ramp shape} answers
the EU-ETS question the conclusions raise.

\paragraph{And it is cheap.} No model change, therefore no recalibration,
therefore no new behavioural assumption to defend. It is the only direction
on this list with essentially zero scientific risk: if the GP fits badly
the diagnostics say so, and the fallback is more design points.

\paragraph{The RL-planner variant.} An alternative to a surrogate over
fixed schedules is a reinforcement-learning planner emitting
$\pi(a_t \mid s_t)$ monthly. This is a different and larger claim --- it
answers ``how much cheaper is 95\% uptake if the regulator can
\emph{adapt}?'' rather than ``what is the best fixed mix?''. It is
attractive because the model is already shaped like an RL environment:
\code{single_policy_simulation} loads a pickled post-calibration
controller and runs only the 144-step future period, so the 64 saved
controllers are ready-made episode resets with built-in seed diversity, and
\code{calc_net_policy_distortion()} is ready-made reward. The main
engineering cost is that \code{manage_policies()} precomputes policy
vectors read by index, so per-step action injection means turning those
reads into a setter. Two design cautions: use common random numbers (fix
the controller batch per update) or gradient variance across 64 stochastic
seeds will dominate; and restrict the observation vector to what a
regulator actually observes (EV share, sales, prices, HHI, fuel prices), not
firm internals, or the learned rule is not implementable and reviewers will
say so.

\paragraph{Recommendation within item 8.} Do the surrogate first --- it is
built, it is the direct answer to a stated limitation, and it produces the
baseline the RL planner would have to beat.

\section{Theme C: supply-side economics}

\subsection{Item 5 --- Financial constraints on firms, budgets for consumers}

\status{not implemented.} The paper's first-named limitation, and it is two
separable ideas.

\paragraph{Firm side.} Section~2.3 states the assumption plainly: firms
``face constant marginal production costs and are not subject to binding
capacity or financial constraints'', defended on the grounds that the
relevant manufacturers are multinationals with capital-market access. The
defence is reasonable and the paper cites Mazzucato and Unruh in
acknowledging what it gives up. Note what the current specification implies:
cumulative firm profit is a \emph{reported output} with no feedback into
innovation capacity. Adding a financing channel closes that loop, which is
attractive because the co-evolutionary feedback \emph{is} the paper's
headline methodological contribution.

\paragraph{Consumer side, which I think is the more interesting half.}
$\beta_i$ scales with the Californian income distribution, but there is no
budget constraint --- a low-income agent can select an arbitrarily expensive
car if the utility arithmetic works out. Figure~6 shows the $\beta_i$
distribution has an extremely long right tail (to $\sim 7\times 10^6$). A
hard budget or a financing constraint (monthly payment as a share of
income, which is how car purchases are actually constrained) would bite
hardest exactly where the paper's used-car-market contribution lives, and
would sharpen the distributional reading of the rebate results. Given that
Figure~7 shows removing the used-car market roughly \emph{triples} EV
uptake by forcing price-sensitive users out of cheap ICEVs, the interaction
between budgets and the used market is likely to be first-order.

\paragraph{Cost.} High. Both halves change calibration-period behaviour, so
recalibration is required, and the consumer-budget version probably requires
re-deriving the analytic $\beta$ median (Supplementary Equation~5), which is
currently pinned by matching the mean observed new-car price.

\subsection{Item 6 --- Switching costs and R\&D costs}

\status{not implemented. The most surgical item in this theme.}

R\&D is currently free: firms are gated only by the activation probability
$\iota$ and by NK locality. Production switching is likewise free, gated
only by $\theta$. So the sole source of inertia in the supply side is
technological path dependence through the landscape --- which is exactly the
``carbon lock-in'' that Section~2.3 cites Unruh (2000, 2002) for and then
argues is not binding.

\paragraph{Why it earns its place.} A per-attempt R\&D cost and a
capital cost of switching a production line give the production subsidy and
the profit channel real teeth, and they give accumulated profit a
\emph{use}. It is also the cheapest way to test whether the paper's
transition speeds are an artefact of costless pivoting. My prior is that
they partly are, and that adding these costs raises the required policy
intensities --- pushing in the same direction as item 4's rationality
channel and the opposite direction to item 1's anticipation channel. The
three together would make a coherent robustness paper on whether the
Table~3 intensities are conservative or optimistic.

\paragraph{Cost.} Moderate. Recalibration required, and one or two new
parameters that will have to be set by assumption plus sensitivity analysis,
in the manner Section~3 already establishes for $\lambda$, $K$ and $\mu$.

\subsection{Item 12 --- Stock, vintage or batch production}

\status{not implemented.}

Production is on demand: Appendix~A.1.1 is explicit that a new design
$y \in Y_t$ ``is represented by a vector of attributes and a sale price, but
does not represent an actual car'', and a car is produced only when bought.
So there is no inventory, no capacity constraint and no lead time.

\paragraph{It is exposed to a headline result.} The paper's second key
takeaway is that ``before 2030, all transition policies entail drawbacks.
The required large-scale deployment of EVs causes an increase in
manufacturing emissions that exceeds the savings in gasoline use.''
Production emissions therefore track sales exactly, month by month, with no
smoothing. Batch or vintage production with inventory would change the
timing and possibly the magnitude of that manufacturing-emissions spike ---
which is the mechanism behind one of the paper's two main policy warnings.
That makes this less of an ``add realism'' item and more of a robustness
check on an existing claim, which is a better framing for it.

\paragraph{Cost.} High, and structural. This is the largest single change to
the model's architecture on the list.

\section{Theme D: technology representation}

\subsection{Item 3 --- Rethinking the NK landscape}

\status{not implemented. Important, but poorly shaped as a paper.}

Two concrete, already-documented problems:

\begin{enumerate}[leftmargin=1.5em]
\item \textbf{Range compression.} Because an attribute's performance is the
  mean of $N=15$ i.i.d.\ draws, normalised performance concentrates in
  $(0.2, 0.8)$ with mean and median $0.5$ (Supplementary Section~3.2). The
  landscape bounds $(F^-, F^+)$ therefore have to be back-solved
  analytically from the target attribute range $(Q^-, Q^+)$ via
  Supplementary Equations~22--28 --- machinery that exists purely to
  compensate for the concentration. The reachable attribute space is
  narrower than the nominal bounds suggest, and gets narrower as $N$ grows,
  which is in tension with wanting $N$ large enough for normality.
\item \textbf{$K$ is confounded with the landscape seed.} The supplementary
  says so itself, at the end of Section~2.1: ``by varying the complexity of
  the landscape, we change the inputs. This means that differences between
  complexity runs may also be due to different seeds in the landscape
  generation.'' So the reported $K$ sensitivity is not clean, and $K_{ICE}$
  is described as having a ``mixed but significant effect'' on EV uptake.
\end{enumerate}

A third, smaller symptom: simulated ICEV driving range is $\sim$20\% too
high (Supplementary Section~1.4), attributed to pinning the fuel tank at
the top tier of observed data.

\paragraph{My view.} Fix the range compression and the $K$-seed confounding
as \emph{infrastructure}, folded into whichever substantive paper comes
next, and report it in that paper's supplementary. Do not try to make it the
paper. There is a publishable version --- ``does the innovation
representation drive the policy conclusion?'', comparing NK against an
alternative such as a continuous cumulative-innovation or learning-curve
formulation --- but it is a methods-venue paper with no policy question
attached, and it would consume a calibration cycle to produce a null result
if the answer is ``no''.

\subsection{Item 13 --- EV battery technology evolution}

\status{partly present as an NK dimension.} Battery capacity is already the
extra fourth EV attribute, bounded $(0, 150)$~kWh from Schmuch et al.\
(2018), with $\rho(C,B)=0$ so that cheap high-capacity EVs are hard to
reach. So this item is largely a \emph{special case of item 3}: making
battery evolution richer means either changing the landscape specification
or bolting an exogenous technology trend onto it. Worth doing as part of
item 3 rather than as a separate direction; the one genuinely distinct
sub-idea is an endogenous learning curve (cost falling in cumulative EV
production), which is a different mechanism from NK search and would
interact with items 6 and 12.

\subsection{Item 14 --- Time-varying carbon intensity of car production}

\status{not implemented. Best cost-to-payoff ratio on the list.}

Manufacturing emissions are \emph{fixed constants} for the whole
2001--2050 horizon: ICEV 10~tonnes, EV 14~tonnes (Table~4, from IPCC 2022).
Battery manufacturing has decarbonised substantially over exactly the period
the model spans, and the model already assumes 90\% grid decarbonisation by
2035 for \emph{driving} emissions while holding \emph{production} emissions
flat --- an internal inconsistency that is easy to state and easy to fix.

\paragraph{Why it matters more than it looks.} It bears directly on two
results. The pre-2030 manufacturing-emissions penalty (Section~4.2.1, and
the third row of Figure~3) is computed with a constant 14~t per EV. And
Supplementary Section~2.4 finds grid decarbonisation matters little for
\emph{uptake} but does matter for emissions \emph{flow}. A declining
production intensity would mechanically improve the EV case and could soften
or shift the ``all policies entail drawbacks before 2030'' conclusion.

\paragraph{Cost.} Low. It is an exogenous time series, structurally
identical to the grid-decarbonisation and fuel-price paths
\code{manage_scenario()} already handles. It does change $\mathrm{LCE}$
during 2001--2023, so it does touch the calibration --- which is an argument
for batching it with item 1 rather than doing it alone.

\section{Theme E: spatial grounding and empirical identification}

\subsection{Items 9 and 10 --- Spatial data and population synthesis}
\label{sec:spatial}

\status{far more built than the idea list suggests.} Commit \code{e7881a9}
(``setting up the code for the zip level synthetic population data set'')
contains a complete, tested-against-placeholder-data pipeline:

\begin{itemize}[leftmargin=1.5em, itemsep=0pt]
\item \code{package/model/} --- \code{synthetic_population.py}
\item \code{package/calibration/} --- \code{summary_stats.py},
  \code{calibration_data_outputs_zip.py},
  \code{NN_multi_round_calibration_zip_gen.py}
\item \code{package/constants/} --- \code{base_params_NN_zip.json}
\item \code{package/generating_data/} --- \code{fake_zip_data_gen.py},
  plus a documented column contract (\code{FAKE_DATA_README.md})
  specifying exactly what real files must supply
\end{itemize}

\noindent
The placeholder population cache (61{,}979 rows) is still in the current
tree under \code{package/calibration_data/FAKE_zip_data/}, carrying
\code{zip}, \code{weight}, \code{income}, \code{vmt}, \code{age},
\code{political_leaning}, \code{ruralness}, \code{latitude} and
\code{longitude}.

The designed summary-statistic vector is 28-dimensional:
\begin{itemize}[leftmargin=1.5em]
\item 8 state EV stock shares (2016--2023) --- the existing aggregate
  anchor;
\item 8 state EV sales shares --- already loaded by the current calibration
  and then never used;
\item 9 cross-sectional gradients: weighted OLS of ZIP EV share on
  z-scored log income, ruralness and political leaning, at three years;
\item 3 between-stratum dispersions of EV share;
\item 3 Moran's~$I$ values of ZIP EV share, on a $k$-nearest-neighbour
  weight matrix built from ZIP centroids.
\end{itemize}

\paragraph{The identification argument, which is the real reason to do this.}
$a_\chi$ is by far the dominant parameter for EV adoption --- first-order
Sobol index $\approx 0.78$ (Figure~S7) --- and it governs the paper's
headline mechanism, the social-imitation threshold. But it is currently
identified from an \emph{aggregate time series} of EV uptake. Aggregate
uptake cannot distinguish social contagion from preference sorting: both
produce S-curves. That is the most obvious available attack on the published
paper, and this pipeline is the answer to it. The design note in
\code{summary_stats.py} states the logic precisely: pure sorting produces a
spatial pattern whose shape is fixed from $t=0$ and merely scales up, while
contagion produces spatial autocorrelation that \emph{grows} as adoption
propagates outward --- so it is the \emph{trajectory} of Moran's~$I$ across
three separated years, not its level, that carries the identifying
information.

It also delivers the paper's own stated route to validation, because
ZIP-level panel data supports genuine training/validation/test splits in a
way one aggregate series does not.

\paragraph{The real blocker is data acquisition, not method.} What is needed
is ZIP-level EV registration stock and sales for California (CEC / DMV
publish at this granularity), ACS income and age by ZCTA, a Census ZCTA
gazetteer for centroids, and a ruralness and political-leaning index. This
is legwork rather than research risk --- which makes it a very different kind
of cost from items 5, 7 or 12.

\paragraph{The one genuine technical blocker.} At $I=3000$ across
$\sim$250 ZIPs, there are roughly 12 agents per ZIP, which is far too noisy
for per-ZIP moments. The gradient design partly works around this by pooling
across ZIPs in a weighted regression rather than binning, but the
population still has to grow substantially --- which is precisely the
``larger car user population to reduce noise \dots through model
parallelisation'' that the conclusions call for. Expect this to be the
dominant engineering cost, and note it compounds with everything else: a
larger population makes every recalibration in every other theme more
expensive.

\subsection{Item 11 --- Charging infrastructure}

\status{not implemented; deliberately scoped out.}

Section~2.2 states that ``the charging infrastructure \dots is a given in
this model'', and Section~3 frames the whole contribution as ``a
transferable framework for evaluating market-based instruments in contexts
where a sufficient charging infrastructure threshold has been reached'',
citing Springel (2021) for the claim that infrastructure has limited
additional effect past a threshold.

\paragraph{This is a stronger scoping argument than the others.} Unlike the
expectations assumption or the constant-production-emissions assumption,
this one comes with a cited empirical justification and an explicit
statement of the domain of validity. Adding infrastructure therefore does
not fix an acknowledged weakness so much as \emph{extend the model to a new
domain} --- jurisdictions below the infrastructure threshold, which is most
of the world and arguably the more policy-relevant case. That is a
legitimate and attractive project, but it should be framed as a domain
extension rather than a correction, and it is only tractable after item 9
because charging supply is intrinsically spatial. Treat it as strictly
downstream of Theme E.

\section{The recalibration bottleneck}
\label{sec:bottleneck}

The single most useful piece of project planning here is to notice which
items share a calibration cycle. An SBI re-run is the expensive, serialising
step, so items that require one should be \emph{batched into a single
recalibration} rather than done sequentially.

\begin{description}[leftmargin=1.5em]
\item[No recalibration:] item 2 (bans are future-period only), item 8
  (operates on saved post-calibration controllers).
\item[Recalibration required:] items 1, 3, 4, 5, 6, 7, 12, 13, 14 (all
  change 2001--2023 behaviour), and items 9, 10, 11 (which redefine the
  calibration target itself).
\end{description}

Note that the two directions requiring no recalibration are exactly the two
that address explicitly stated limitations. That is not a coincidence: the
paper scoped out what was expensive.

A corollary for Theme E: because items 9 and 10 \emph{change the calibration
target} rather than merely invalidating the posterior, they should come
either first --- before behavioural changes accumulate --- or last, after
they have all settled. Interleaving a spatial recalibration with behavioural
extensions means paying for both at once and being unable to attribute a
change in fit to either.

\section{Summary}
\label{sec:table}

\begin{longtable}{@{}p{0.20\textwidth}cccp{0.30\textwidth}@{}}
\toprule
\textbf{Item} & \textbf{Built?} & \textbf{Recal.} & \textbf{Cost} &
\textbf{What it buys} \\
\midrule
\endhead
8 Multi-policy optimisation & mostly & no & low &
Full 5-instrument Pareto front; time-varying schedules. Lifts a stated
limitation. \\
\addlinespace
2 Command and control & mostly & no & low &
ZEV mandate on the ZEV state. Needs item 1. \\
\addlinespace
1 Forward-looking agents & yes & yes & low &
Removes the permanent-policy assumption. Effect known to be large. \\
\addlinespace
14 Production carbon intensity & no & yes & low &
Fixes an internal inconsistency; touches the pre-2030 emissions result. \\
\addlinespace
4 Long-horizon firm profit & no & yes & low--mod &
Aligns the code with what Section~2.3 already claims. Non-RL route to most
of item 7. \\
\addlinespace
6 Switching / R\&D costs & no & yes & mod &
Tests whether transition speeds are an artefact of costless pivoting. \\
\addlinespace
3 NK landscape & no & yes & mod &
Fixes range compression and $K$-seed confounding. Infrastructure, not a
paper. \\
\addlinespace
13 Battery evolution & partly & yes & mod &
Largely a special case of item 3, except the learning-curve variant. \\
\addlinespace
9 + 10 Spatial / pop.\ synthesis & mostly & target &
high & Identifies contagion vs.\ sorting. Defends the core mechanism.
Blocker is data, not method. \\
\addlinespace
5 Financial / budget constraints & no & yes & high &
Paper's first-named limitation. Consumer-budget half interacts with the
used-car contribution. \\
\addlinespace
12 Batch / vintage production & no & yes & high &
Robustness on the manufacturing-emissions spike. Largest architectural
change. \\
\addlinespace
11 Charging infrastructure & no & yes & high &
Extends the model below the infrastructure threshold. Downstream of item 9.
\\
\addlinespace
7 MARL firms & no & yes & very high &
Commitment vs.\ discretion. Presupposes item 8; breaks the calibration
chain. \\
\bottomrule
\end{longtable}

\section{Recommended sequence}

\paragraph{First, and immediately: item 8 (surrogate).}
It is built, it needs no recalibration, it introduces no behavioural
assumption to defend, and it answers a question the paper explicitly says it
could not afford to answer. Restore \code{package/surrogate/} from
\code{c92bff8}, re-point it at the current
\code{base_params_endogenous_policy_pair_gen.json}, validate the GP
($R^2$, leave-one-out, interval coverage), and extend the design space from
five intensities to five intensities plus ramp shape. Lowest-risk paper
available.

\paragraph{Second: items 1 + 2 + 4 as one bundle --- ``anticipation and
regulation''.}
One recalibration buys all three. Item 1 is already implemented; item 2's
three bans are already implemented; item 4 is a contained change to one
profit calculation. The motivation is quoted verbatim from the paper's own
conclusions, the case study is the jurisdiction that actually has a ZEV
mandate, and the driving-ban smoke test already indicates the effect is
large rather than marginal. Fold item 14 into the same recalibration since
it is nearly free once you are paying for one.

\paragraph{Third, as the long-term bet: Theme E (items 9 + 10).}
This is the one that protects the published paper. $a_\chi$ carries the
headline mechanism and dominates the Sobol decomposition, and it is
currently identified from an aggregate series that cannot separate contagion
from sorting. The pipeline for fixing that is designed, written and tested
against placeholder data; what is missing is real ZIP-level California data
and a larger population. Start the data acquisition now, in parallel with
the first two projects, because it is the long-lead item and it is
administrative rather than intellectual work.

\paragraph{Defer: items 7, 12, 5.} Each is defensible and each is
expensive. Item 7 in particular should wait until item 8's RL-planner
variant exists, because without an adaptive planner there is nothing for
learning firms to learn that reading the announced schedule would not
deliver. Items 5 and 12 are natural if a collaborator wants the industrial
-organisation angle; they are not the best use of a solo calibration cycle.

\paragraph{Fold in rather than pursue: items 3, 13.} Fix the NK range
compression and the $K$-seed confounding inside whichever project comes
next and report it in that supplementary.

\paragraph{A caution on breadth.} Fourteen directions is a portfolio, not a
plan, and they point in genuinely different directions: Theme B is
computational, Theme A is behavioural, Theme C is industrial organisation,
Theme E is econometric. Three of these done properly is a stronger position
than eight touched. The sequence above is chosen so that each project
leaves behind infrastructure the next one needs --- the surrogate becomes
the RL planner's baseline, the anticipation bundle becomes the environment
strategic firms would learn in, and the spatial calibration becomes the
identification backbone for everything after.

\section{One direction the paper promises that is not on this list}
\label{sec:missing}

The conclusions commit to expanding the social-learning representation
``beyond a threshold model to enrich car users with social norms and
behavioural control such as in the Theory of Planned Behaviour (Ajzen,
1991)'', explicitly in order to ``facilitate the inclusion of policies such
as information provision, media or marketing''.

This is worth adding to the portfolio for three reasons. It is the demand
-side counterpart to Theme A --- everything in that theme makes agents
smarter about \emph{prices}, and nothing makes them richer about
\emph{norms}. It unlocks a whole policy class (information, media,
marketing) that the current model cannot represent at all, complementing
the market-based instruments already characterised and the regulatory ones
item 2 would add. And it bears on the same identification problem as
Theme E: the binary consideration threshold $x_{i,t}$ is a very sharp
behavioural assumption carrying a great deal of the model's dynamics, and a
graded alternative is both more realistic and a natural robustness check on
whether the social-lock-in non-linearity of Figure~S13 survives a softer
specification.

\end{document}
